% sample.tex -- GoodCompanions2 Chess Font
%   Gallery page in the style of chessfonts_gallery.pdf (Ulrike Fischer)
% Compile with: pdflatex sample.tex
% (Font must be installed -- see README.md for installation instructions.)
\documentclass{article}

\usepackage[
  a4paper,
  top=2.5cm, bottom=2.5cm, left=3cm, right=3cm
]{geometry}

\usepackage[LSB1,LSB2,LSB3,LSF,T1]{fontenc}
\usepackage{xcolor}
\usepackage{array}
\usepackage{booktabs}
\usepackage{chessboard}
\usepackage{chessfss}

% ---- Font setup ----
\setfigfontfamily{goodcompanions2}

\setchessboard{
  boardfontfamily=goodcompanions2,
  boardfontencoding=LSB1,
  boardfontsize=20pt,
  borderwidth=0.5pt,
  bordercolor=black,
  showmover=false,
  labelleft=true,
  labelbottom=true,
  labelfont=\rmfamily\bfseries,
  labelfontsize=7pt
}

\newcommand{\startfen}{rnbqkbnr/pppppppp/8/3qk3/3QK3/8/PPPPPPPP/RNBQKBNR w KQkq - 0 1}

\begin{document}
\parindent=0pt

\bigskip

\noindent\textbf{GoodCompanions2}

\smallskip
\begin{tabular}{@{}ll@{}}
source:     & \texttt{https://github.com/joelcato/goodcompanions2-chess-font} \\[2pt]
author:     & Joel Cato \\[2pt]
based on:   & GoodCompanions by David L.\ Brown (2004) \\[2pt]
characters: & figurine and board symbols \\[2pt]
familyname: & goodcompanions2 \\[2pt]
fonts:      & The following fonts for the package chessfss can be made from this source: \\
\end{tabular}

\medskip
\begin{tabular}{@{}lll>{\raggedright\arraybackslash\small}p{6.5cm}@{}}
\toprule
encoding & serie & tfm-name & reencoding command for \texttt{chess.map} \\
\midrule
raw & --          & \texttt{chess-goodcompanions2-board-fig-raw} & none \\[4pt]
LSF & \texttt{m}  & \texttt{chess-goodcompanions2-lsf}
    & \texttt{" ChessFigEncoding ReEncodeFont "}
      \texttt{<chess-fig.enc} \\[4pt]
LSB & \texttt{m}  & \texttt{chess-goodcompanions2-lsb}
    & \texttt{" ChessBoardEncoding ReEncodeFont "}
      \texttt{<chess-board.enc} \\
\bottomrule
\end{tabular}

\bigskip

% ---- Four chessboards (2x2 grid) ----

\noindent
\begin{minipage}[t]{0.47\textwidth}
  \chessboard[setfen=\startfen,
    clearfontcolors,blackfieldcolor=black!60,setfontcolors]
\end{minipage}\hfill
\begin{minipage}[t]{0.47\textwidth}
  \chessboard[setfen=\startfen,
    clearfontcolors,blackfieldcolor=black!60,colorwhite=blue,colorblack=red,setfontcolors]
\end{minipage}

\bigskip

\noindent
\begin{minipage}[t]{0.47\textwidth}
  {\setchessboard{boardfontencoding=LSB2}%
  \chessboard[setfen=\startfen,
    color=black!30,colorblackbackfields,
    clearfontcolors,colorwhite=blue,colorblack=red,setfontcolors]}%
\end{minipage}\hfill
\begin{minipage}[t]{0.47\textwidth}
  {\setchessboard{boardfontencoding=LSB3}%
  \chessboard[setfen=\startfen,
    color=black!30,colorblackbackfields,
    clearfontcolors,colorwhite=blue,colorblack=red,setfontcolors]}%
\end{minipage}

\bigskip

% ---- Figurine specimen (LSF encoding) ----

{\large\figfont\symking\symqueen\symrook\symbishop\symknight\sympawn}\\[4pt]

\newpage

% =============================================================
%  Merida
% =============================================================

\setfigfontfamily{merida}
\setchessboard{boardfontfamily=merida}

\noindent\textbf{Merida}

\smallskip
\begin{tabular}{@{}ll@{}}
source:     & \texttt{https://www.enpassant.dk/chess/fonteng.htm} \\[2pt]
author:     & Armando H.\ Marroquin (1998) \\[2pt]
characters: & figurine and board symbols \\[2pt]
familyname: & merida \\[2pt]
fonts:      & The following fonts for the package chessfss can be made from this source: \\
\end{tabular}

\medskip
\begin{tabular}{@{}lll>{\raggedright\arraybackslash\small}p{6.5cm}@{}}
\toprule
encoding & serie & tfm-name & reencoding command for \texttt{chess.map} \\
\midrule
raw & --          & \texttt{chess-merida-board-fig-raw} & none \\[4pt]
LSF & \texttt{m}  & \texttt{chess-merida-lsf}
    & \texttt{" ChessFigEncoding ReEncodeFont "}
      \texttt{<chess-fig.enc} \\[4pt]
LSB & \texttt{m}  & \texttt{chess-merida-lsb}
    & \texttt{" ChessBoardEncoding ReEncodeFont "}
      \texttt{<chess-board.enc} \\
\bottomrule
\end{tabular}

\bigskip

\noindent
\begin{minipage}[t]{0.47\textwidth}
  \chessboard[setfen=\startfen,
    clearfontcolors,blackfieldcolor=black!60,setfontcolors]
\end{minipage}\hfill
\begin{minipage}[t]{0.47\textwidth}
  \chessboard[setfen=\startfen,
    clearfontcolors,blackfieldcolor=black!60,colorwhite=blue,colorblack=red,setfontcolors]
\end{minipage}

\bigskip

\noindent
\begin{minipage}[t]{0.47\textwidth}
  {\setchessboard{boardfontencoding=LSB2}%
  \chessboard[setfen=\startfen,
    color=black!30,colorblackbackfields,
    clearfontcolors,colorwhite=blue,colorblack=red,setfontcolors]}%
\end{minipage}\hfill
\begin{minipage}[t]{0.47\textwidth}
  {\setchessboard{boardfontencoding=LSB3}%
  \chessboard[setfen=\startfen,
    color=black!30,colorblackbackfields,
    clearfontcolors,colorwhite=blue,colorblack=red,setfontcolors]}%
\end{minipage}

\bigskip

{\large\figfont\symking\symqueen\symrook\symbishop\symknight\sympawn}\\[4pt]

\newpage

% =============================================================
%  SKAKnew
% =============================================================

\setfigfontfamily{skaknew}
\setchessboard{boardfontfamily=skaknew}

\noindent\textbf{SKAKnew}

\smallskip
\begin{tabular}{@{}ll@{}}
source:     & \texttt{https://www.enpassant.dk/chess/fonteng.htm} \\[2pt]
author:     & Eric Bentzen (2003) \\[2pt]
characters: & figurine and board symbols \\[2pt]
familyname: & skaknew \\[2pt]
fonts:      & The following fonts for the package chessfss can be made from this source: \\
\end{tabular}

\medskip
\begin{tabular}{@{}lll>{\raggedright\arraybackslash\small}p{6.5cm}@{}}
\toprule
encoding & serie & tfm-name & reencoding command for \texttt{chess.map} \\
\midrule
LSB (T) & \texttt{m}  & \texttt{SkakNew-DiagramT} & (sizes $<$16pt) \\[4pt]
LSB     & \texttt{m}  & \texttt{SkakNew-Diagram}  & (sizes $\geq$16pt) \\[4pt]
LSF     & \texttt{m}  & \texttt{SkakNew-Figurine} & \\
\bottomrule
\end{tabular}

\bigskip

\noindent
\begin{minipage}[t]{0.47\textwidth}
  \chessboard[setfen=\startfen,
    clearfontcolors,blackfieldcolor=black!60,setfontcolors]
\end{minipage}\hfill
\begin{minipage}[t]{0.47\textwidth}
  \chessboard[setfen=\startfen,
    clearfontcolors,blackfieldcolor=black!60,colorwhite=blue,colorblack=red,setfontcolors]
\end{minipage}

\bigskip

{\large\figfont\symking\symqueen\symrook\symbishop\symknight\sympawn}\\[4pt]

\end{document}
